\gddchapter{Interview sequence}


Below is the guideline for the semi-open guided interview. 
The interview starts with easy warm-up questions and progresses towards harder questions as it goes along. 

The answers were translated from polish back to english.

\section{Introduction}
Remind the participant that this is not a test, rather a joint exploration.

\section{Icebreakers}

\begin{QandA}
   \item Now that you’ve tried both versions, how are you feeling? Was there anything that immediately surprised you?
         \begin{answered}
            Maria: The second one was better.

            I had choices instead of speaking.
            I preferred either full interaction by speaking like in a chat with AI, or using button commands.
            understood the game better with commands.

            Gabriel: So the second version.

            Gabriel: Was there anything surprising at the beginning?

            Maria: I felt overwhelmed and didn’t know what to do.
            I didn’t know how to start.
            That was the hardest part.
             But I also don’t have experience with such games.

            \end{answered}

   \item If you had to describe the difference between these two games to a person who hasn't played either, what would you say?
         \begin{answered}
               Maria: The difference was that pressing keys felt more intuitive than speaking to the computer.

               Speaking felt unnatural because I’m not used to telling a game what I want to do.

               Gabriel: You mentioned it was overwhelming—can you explain more how using voice felt?

               Maria: I felt a language barrier because I’m not as comfortable in English as I am with reading and pressing keys.
               When I read text and commands, I fully understand them.
               But speaking was much harder and more tiring.
               If it were in Polish, it would be much more enjoyable.
               If it were in a language I’m fully comfortable with, I would try again.

         \end{answered}
         
\end{QandA}




\section{Grand tour questions}
(remember to pause here 3-5 seconds)

\begin{QandA}


   \item From the moment you first used your voice to navigate, what were you thinking as you decided what to say to the system?

         \begin{answered}
            No, I was basically using the simplest words because I know it’s an AI.
            In the sense that I do it so that it understands it.
            I identify with the character, that it’s me.
            And I’m not saying “he goes there” and “I’m going there.”
         \end{answered}

    \item How was it like to use the voice input in the game?
         \begin{answered}
            Maria: The difference was that pressing keys felt more intuitive than speaking to the computer.
            felt unnatural because I’m not used to telling a game what I want to do.

            Gabriel: You mentioned it was overwhelming—can you explain more how using voice felt?

            Maria: I felt a language barrier because I’m not as comfortable in English as I am with reading and pressing keys.
            When I read text and commands, I fully understand them.
            But speaking was much harder and more tiring.
            If it were in Polish, it would be much more enjoyable.
            If it were in a language I’m fully comfortable with, I would try again.

         \end{answered}


\end{QandA}




\section{Core comparative questions}

\begin{QandA}
   \item In which version did you feel more 'inside' the story world?
         \begin{answered}
            Maria: The first one.
            Because I didn’t have to read text—I just looked at the image and imagined the world.
            I would prefer the first version if I didn’t have to read at all.
            If I spoke, the computer could read the text to me.
            Even in a mechanical voice.
            That would be very valuable.
            So it would read it to you?
            Yes, it would read it to me.
            And during the game, it would also guide me and ask questions like:
            “What do you want to do now?”
            It would guide me a bit.

            Gabriel: OK, something like that.

            Maria: Yes.
            What are we doing now?
            Then I would have choices and could respond.
            I would also like to express emotional information—what I feel when entering a place.
            Like “I feel afraid to be here” or “I have a bad feeling.”
            It’s like in a movie when a character enters a room and senses something bad will happen.
            Or something interesting.
            I’d like more freedom of movement and expression.
            Then I’d prefer that over selecting options, because options are limiting.
            Even speaking would let me stay longer in a place and explore more.

         \end{answered}

   \item How did the 'effort' of thinking of what to say naturally compare to the 'effort' of knowing which keys to press?

         \begin{answered}
         Maria: It’s definitely more effort with voice commands.
         With keys it’s already there—I just press.

         Gabriel: I understand.

         Maria: That’s relaxing.
         With voice, I have to think about what to say.
         I don’t have enough data to feel comfortable.

         \end{answered}

    \item Did the freedom to say 'anything' make you feel more in control, or was it more overwhelming than having a fixed list of keys?

         \begin{answered}
            It was more uncertain.
            I didn’t know what I should do.
            I felt lost..
         \end{answered}
    \item How would you compare the overall experience of using the two game versions? Which one would you choose and why?

         \begin{answered}
            Maria: If nothing changed right now?

            Gabriel: Yes.

            Maria: I would choose the keyboard version because I understood it faster.

            Gabriel: OK.

            Maria: Because I had commands like “use item,” so I knew what to do.

            I could go back and search rooms.

            Gabriel: OK.

            Maria: That was better for me.

            Gabriel: OK.

            Gabriel: And comparing the experiences overall?

            Maria: The first version was more strange and unfamiliar, while the second was more familiar because I’ve used keyboard controls in games before.
            It felt more intuitive.

            Gabriel: OK.

            At one point you asked if you can ask the game questions.

            Maria: Yes.

            Gabriel: What kind of questions?

            Maria: What should I do now?
            Or what’s behind these doors?
            Something that expands the world.

            Can I use a crowbar?
            Is there another floor?
            Questions that expand the map.

            Gabriel: And when you used voice commands, what guided your choice of words?

            Maria: Words I already knew.
            I imagined keyboard commands—go forward, run, crouch, stay, go back.
            Movement-based commands.
            
         \end{answered}
\end{QandA}






\section{Probing}
In this part of the interview look into the sticky moments noted above and ask follow-up questions probing for more detail.
Include lexical reflection (eg. you used the word "clunky" when explaining navigation. Could you tell me more about what was happening
in that moment?)

The probing was done as follow up questions to the pre-existing questions and in the discussion during the last question rather than separate questions in this test.

% \begin{QandA}
%    \item It looks like you really value creativity when it comes to playing games. Would you agree with that?
%    \begin{answered}
%       I would agree with that because
%       If I played, let's say, a lot of games, then they can become kind of dull.
%       So if something is new for me, or it's like, you know, creative or interesting,
%       I value that because I'm pretty worn out by some normal or usual type of games like the second one.
%       It was way more similar to the games that I played in the past, I would say.
%    \end{answered}
% \end{QandA}


\section{Final reflection}

\begin{QandA}
   \item After reflecting on everything today, is there anything else you’d like to add about using your voice to play this game that we haven't talked about?

         \begin{answered}
            Maria: I would like more connections between rooms.
            And I would like my own music.

            Gabriel: Audio?

            Maria: Yes.

            Maria: Probably.

            Gabriel: What do you mean by connections between rooms?

            Maria: For example, earlier I couldn’t move from storage to another part of the mall.

            Gabriel: So not going from A to C directly, but through B?

            Maria: Yes.

            Gabriel: OK.

            Maria: I’d like more exploration like hidden paths and shortcuts.

            I wanted to use the crowbar in unusual ways, like breaking a cabinet in the Korean store to discover something.
            It made me feel a bit aggressive.
            I wanted to interact more destructively.

            Gabriel: I understand.

            Maria: There were also no real threats.
            Like maybe a dangerous dog in a room.

            Gabriel: I understand.

            Maria: Something I had to deal with.

            Gabriel: Got it.

            Maria: But I’m not sure I could handle that since I struggled already.
         \end{answered}
\end{QandA}

\section{Closing}
Thank the participant for his time and tell them again how much value his input and time gives me.






